\documentclass[12pt,a4paper]{article}
% ---- Basic packages ----
\usepackage[utf8]{inputenc}
\usepackage[T1]{fontenc}
\usepackage[english]{babel}
\usepackage{lmodern}
% ---- Page layout ----
\usepackage{geometry}
\geometry{a4paper, margin=2.5cm}
\usepackage{enumitem}
% ---- Line spacing ----
\usepackage{setspace}
\onehalfspacing
% ---- Math and symbols ----
\usepackage{amsmath, amssymb, amsfonts}
\usepackage{bm}
\usepackage{siunitx}
% ---- Tables ----
\usepackage{booktabs, tabularx, longtable, multirow, threeparttable}
% ---- Figures ----
\usepackage{graphicx}
\usepackage{float}
\usepackage{caption}
\usepackage{subcaption}
% ---- Colors and links ----
\usepackage{xcolor}
\usepackage{hyperref}


\title{Geopolitics for Business \\ Group Project}

\author{
    Data Group 2 \\[2ex]
    \normalsize Marius Borsu Lefevre \\
    \normalsize Giovanni Clemente \\
    \normalsize Greta Guadalupi \\
    \normalsize Petr Hrebenar \\
    \normalsize Stéphane Miquel Clément \\
    \normalsize Pietro Montano \\
    \normalsize Souhail Habib Nehme \\
    \normalsize Joan Olaf Subirachs Garcia
}

\date{\today}

\begin{document}
\maketitle
\newpage
\tableofcontents
\newpage

% =============================================================================
\section{Introduction and Research Question}
% =============================================================================

Russia's full-scale invasion of Ukraine in February 2022 constituted one of the most significant energy supply disruptions in recent European history. The sanctions imposed on Russia and the subsequent re-routing of European energy procurement transformed what had been a decade-long exposure to cheap Russian natural gas into an acute strategic liability, particularly for countries and industries that had built their production models around it. At the same time, the shock arrived against a backdrop of intensifying climate commitments: the EU's Fit for 55 package and the corporate disclosure obligations under the Corporate Sustainability Reporting Directive (CSRD) were already pushing firms to reduce their greenhouse gas emissions per unit of output.

This paper asks how the Russia–Ukraine energy shock affected firm-level greenhouse gas intensity. Specifically, we are interested in whether European firms — and in particular, energy-intensive manufacturers — reduced their direct GHG emissions per unit of revenue in the aftermath of the shock, and whether this adjustment was shaped by the degree to which their host country was exposed to Russian gas imports. The core research question is therefore: \textit{Did the Russia–Ukraine war trigger a reduction in firm-level GHG intensity among European energy-intensive firms, and was this effect attenuated for firms operating in countries that were more dependent on Russian gas?}

The question sits at the intersection of two literatures that have thus far developed in relative isolation: the economics of energy market disruptions and the firm-level determinants of environmental performance. It also carries direct policy implications. If energy-intensive firms in gas-dependent countries were less able to reduce their GHG footprint — because the menu of available energy alternatives was more constrained — then country-level energy policy becomes a key enabler (or bottleneck) for corporate decarbonisation. The analysis proceeds as follows. Section~\ref{sec:lit} situates the paper in the relevant literature. Section~\ref{sec:data} describes the data construction and treatment assignment. Section~\ref{sec:method} presents the empirical strategy. Section~\ref{sec:results} discusses results and their interpretation. Section~\ref{sec:limits} outlines limitations and avenues for future research. Section~\ref{sec:conclusion} concludes.


% =============================================================================
\section{Literature Review}
\label{sec:lit}
% =============================================================================

\subsection{Theoretical Framework}

Academic and professional notions of the condition of globality, widespread first during the 1980s (Robertson, 1983), emerged along with an increasing cross-border integration of the global economic system resulting after the II World War (Hover, Buchenrieder, and Sokolsky, 2025). Indeed, the characterisation of a liberal regime of financial and trade relations based on the principles of cooperation, mutual interdependence, and openness, prompted the flourishing of multinational enterprises (MNEs) in the international business (IB) environment (Petricevic and Teece 2019; Lessard, Teece, and Leih, 2016). The liberal school of thought, acknowledging the increasing ``flatness'' of the system (Friedman, 2005; Rugman and Oh, 2008), argued that the resulting deepening of economic integration would diminish the likelihood of inter-state confrontation due to high opportunity costs from the disruption of trade (Mansfield and Pollins, 2003; Choi and Kim, 2021; Gartzke and Li, 2003; Goldsmith, 2013). This political perspective was operationalised, for instance, through Germany's ``change through trade'' approach that encouraged the construction of several gas pipelines with the Soviet Union and later Russian Federation between the 1970s and the early 2000s (Blumenau, 2022; Wood, 2023). In light of this, energy infrastructure projects such as Nord Stream 2 were ``never just business'' (Schmidt-Felzmann, 2016). However, geopolitical developments over the past fifteen years have revealed the limitations of the liberal theory, indicating a reversal of an assumption that is reshaping the business landscape once again.

In this context, several authors have put forward that the current world order is entering a de-globalisation phase (Livesay, 2018; James, 2018); one that is putting at risk, by direct means, the way MNEs operate and do business (Petricevic and Teece, 2019). More specifically, this period of disintegration is making the IB landscape fraught with unprecedented levels of volatility, uncertainty, complexity, and ambiguity (VUCA) (Bennett and Lemoine, 2014; Cavusgil et al., 2021; Schoemaker, Heaton, and Teece, 2018; van Tulder, Verbeke and Jankowska, 2020). As established by Sénquiz-Díaz (2025), the most important contribution to the VUCA framework in the IB context is that of Bennett and Lemoine (2014). For these authors, volatility can be described as a state of ``relatively unstable change'', while uncertainty is the ``lack of knowledge as to whether an event will have meaningful ramifications'' (p.\ 313). Furthermore, complexity, in the IB environment, reflects the idea that there is no single arrangement, but rather a constellation of systems that interact altogether in an exhaustive ``network of information of procedures'' (p.\ 313). Lastly, ambiguity is depicted as the absence of clear guidelines for decision-making, and the margin for independent interpretation in cause-and-effect relationships.

MNEs are currently operating in such one highly challenging context that not only disrupts routine business operations but also complicates the management of exceptional circumstances. As put forward by Cavusgil et al.\ (2021), apart from traditional commercial, financial, cultural, and political risks, companies are now dealing with additional and greater threats, including technological obsolescence, cyber security issues, trade wars, sustainability concerns, social challenges, and, more importantly, ``black swan'' events. Indeed, in the systems risk literature, exogenous shocks such as interstate armed conflicts or extreme weather events emerge as ``crises'' once actual harms are materialised (Homer-Dixon et al., 2015). Nevertheless, as put forward by Lawrence, Janzwood, and Homer-Dixon (2022), today, the causal activation of a given risk in a given system into an actual crisis does not remain confined to the system in which it originates, but it spreads and reinvigorates onto other systems.

The term \textit{polycrisis} reflects the confluence and intertwining of these individual crises that ``interact, exacerbate, and reshape one another (\ldots) that must be understood and addressed as a whole'' (Lawrence, Janzwood, and Homer-Dixon, 2022). The conceptual transition from systemic risk to the notion of polycrisis can be understood through the maximisation of the broader defining properties embedded in the systems that the latter emphasise. Based on the scholarship analysis of Lawrence, Janzwood, and Homer-Dixon (2022), the five distinct characteristics of polycrises are: extreme complexity and dynamism; disproportional cause-effect relationships, transboundary causal links, and profound uncertainty (Renn et al., 2020; 2019; Schweizer and Renn, 2019). Precisely, as put forward by Collste et al.\ (2025), the identification of valid causality relationships between these crises under such uncertain conditions makes it hard to assess and overcome the challenges derived from it. In the application of a system archetypes approach to identify causal trends within crises, the authors define a model, referred to as ``limits to growth'', that effectively specifies our scope of research. In the context of a growing size of economic operations, certain operation constraints and external pressures such as scarce resources or climate-related externalities constrain the expansion capacity of operations of businesses. Furthermore, in light of this increased complexity, Rapoza (2019) argues that MNEs might be already shifting toward decoupling from their global and regional global value chains (GVCs), a strategy also envisioned by Collste et al.\ (2025) to manage the implications of polycrises.

This strategic decoupling from global and regional value chains cannot be fully understood without first examining the theoretical architecture that has governed the organisation of international production over the past decades. Global Value Chain (GVC) theory and Global Production Network (GPN) theory provide complementary analytical lenses through which the structural vulnerabilities and strategic reconfiguration of MNEs can be assessed in light of compounding exogenous shocks. The foundational contribution to GVC governance was advanced by Gereffi, Humphrey, and Sturgeon (2005), who identified five distinct governance types ranging from low to high levels of explicit coordination and power asymmetry, determined by transaction complexity, codifiability, and supply-base capabilities. In the context of the VUCA landscape described above, these determinants are precisely the variables most susceptible to disruption: geopolitical fragmentation raises transaction complexity, energy and trade shocks erode the codifiability of cross-border exchanges, and climate-related pressures strain supplier capabilities across multiple geographies simultaneously. Yeung and Coe (2015) advanced GPN 2.0, a more explicitly causal and dynamic theory designed to explain why and how the organisation and coordination of global production networks varies significantly within and across industries, sectors, and economies. The relational and power dimensions of these networks further illuminate how shocks propagate and are redistributed across chain actors (Kano, 2018; Humphrey and Schmitz, 2002). Ponte (2020) further demonstrates that lead firms' pursuit of environmental upgrading frequently results in the upstream displacement of sustainability compliance costs and risks onto suppliers, a finding directly pertinent to our research question: climate shocks do not affect GVC actors uniformly, but interact with pre-existing governance asymmetries to produce differentiated strategic and financial outcomes across the network.

\subsection{Empirical Approach and Choice of Metrics}

A growing literature shows that geopolitical shocks affect firms not only through standard financial variables such as investment, financing costs, and profitability, but also through sustainability-related outcomes. In particular, recent firm-level studies find that higher geopolitical risk and geopolitical conflict are associated with weaker ESG performance, including weaker environmental outcomes (Abdullah et al., 2024; Saharti et al., 2024; Erzurumlu et al., 2025). This means that environmental performance is not a peripheral variable in this setting, but one of the firm-level margins through which geopolitical stress can be observed.

Methodologically, difference-in-differences (DiD) has become a widely used approach for identifying causal effects of shocks and policy changes. Fan et al. use the 2018 US–China trade war as a quasi-exogenous shock and compare U.S. firms with direct Chinese suppliers to firms without such suppliers in a DiD framework, focusing on operating-performance outcomes such as inventory days and return on assets. 
More relevant to our specification is the triple-difference extension: Fan et al.\ (2022) also estimate a DDD to assess whether the trade war effect on US firms was moderated by the degree of supply-chain diversification, exploiting variation across firms within treated and control groups. As noted by Olden and Moen (2022), the parallel trends assumption in a DDD holds if the difference between treated and control subjects in a treated state equals the corresponding difference in a control state, a weaker requirement that our triple-interaction Model (4) is designed to exploit.

Once environmental outcomes are accepted as relevant, the next issue is how to measure them. Total capital expenditure (CAPEX) and CSR spending are both informative, but neither is a direct measure of environmental performance: total CAPEX does not distinguish between investment that reduces environmental harm and investment that maintains or expands emissions-intensive activity, while CSR expenditure may reflect communication or stakeholder management rather than operational performance. For that reason, an environmental metric tied directly to emissions is closer to the underlying object of interest (Abdullah et al., 2024). This is consistent with the broader climate-finance literature, where carbon emissions and carbon intensity are regularly used to capture firms' exposure to climate-related risks and their environmental footprint (Sautner et al., 2023; Bolton and Kacperczyk, 2023; Perdichizzi et al., 2024).

Within that class of metrics, Trucost's GHG Direct Intensity captures greenhouse-gas emissions generated by the firm's own operations (Scope 1) relative to revenue. This makes it an operational environmental-performance variable: it measures how emissions-intensive a firm's own production is, rather than how exposed it is to controversies or how strong its sustainability communication appears to be (S\&P Global, 2025; EPA, 2025). RepRisk was considered as an alternative metric, but it measures externally reported environmental controversies rather than firms' underlying operational environmental efficiency; its methodology is explicitly outside-in and issues-driven, making it suitable for tracking incident exposure but less directly tied to emissions intensity. In addition, Barkemeyer, Revelli, and Douaud (2023) show that the odds of controversy coverage are five times higher for firms headquartered in English-language countries, introducing a systematic media-selection bias. For that reason, RepRisk is better interpreted as a controversy measure than as a measure of operational environmental performance. There is also a measurement reason to focus on direct emissions: the NGFS (2024) finds that inconsistencies across emissions datasets are lowest for Scope 1, while discrepancies increase substantially for Scope 2 and especially Scope 3. The main limitation of Trucost GHG Direct Intensity is frequency, assessments are aligned with firms' annual reporting cycles, but the main studies linking geopolitical shocks to ESG outcomes also rely on annual panels (Abdullah et al., 2024; Saharti et al., 2024; Erzurumlu et al., 2025), so the use of an annual environmental variable is consistent with the frequency that already dominates this literature.

Building on this theoretical and empirical foundation, our study focuses on the Russia--Ukraine energy shock as a concrete instance of polycrisis dynamics, exploiting cross-country variation in Russian gas dependency and cross-sector variation in energy intensity to identify its effect on firm-level GHG performance.


% =============================================================================
\section{Data}
\label{sec:data}
% =============================================================================

\subsection{Overview of the Panel Dataset}

The final dataset is a firm-year panel covering publicly listed European firms in selected manufacturing and resource sectors. It combines three data sources: proprietary ESG firm-level data from S\&P Global TruCost; country-year macroeconomic data from the World Bank; and country-year energy price data from EMAPS/Ember. The panel spans the years 2017–2024 and covers 28 European countries (EU-27 plus Norway, Switzerland, and the United Kingdom). After filtering for the relevant sector universe, removing duplicate firm-year observations, and excluding singleton firms (which contribute no within-variation under a firm fixed effects estimator), the final panel contains approximately 973,000 firm-year observations across roughly 192,000 unique firms.

\subsection{Outcome Variable: GHG Direct Intensity}

The primary outcome variable is \textit{GHG Direct Intensity}, defined as Scope 1 (direct) greenhouse gas emissions in tonnes of CO$_2$-equivalent per USD 1 million of revenue, as reported by TruCost. As argued in detail in Section~\ref{sec:lit}, this measure is the most appropriate for our setting: it captures the operational emissions efficiency of the firm, is directly linked to energy input decisions, normalises for firm size, and benefits from the lowest measurement noise across emissions categories (NGFS, 2024). Among the alternatives considered — total CAPEX, CSR spending, and the RepRisk controversy index — RepRisk was the closest substitute in concept, but proved unsuitable due to substantial coverage gaps across our European sample and a systematic media-selection bias that disproportionately reflects firms domiciled in English-language countries (Barkemeyer, Revelli, and Douaud, 2023). We therefore proceed with GHG Direct Intensity, which has near-complete coverage across the full panel.

Table~\ref{tab:desc} presents descriptive statistics for the key variables.

\begin{table}[H]
\centering
\small
\caption{Descriptive Statistics}
\label{tab:desc}
\resizebox{\textwidth}{!}{%
\begin{tabular}{lrrrrrrr}
\toprule
Variable & N & Mean & SD & p25 & Median & p75 & Max \\
\midrule
GHG Direct Intensity (tCO$_2$e / \$M rev.) & 973,815 & 121.1 & 205.3 & 16.5 & 88.7 & 188.4 & 15,225 \\
GHG Scope 1 Absolute (tCO$_2$e) & 973,815 & 17,081 & 761,962 & 30 & 157 & 657 & 181,534,000 \\
Revenue (USD million) & 973,815 & 94.1 & 2,424 & 0.6 & 2.4 & 10.1 & 400,223 \\
GDP per Capita (EUR) & 973,445 & 48,141 & 12,768 & 40,353 & 46,680 & 54,688 & 140,739 \\
Inflation (\%) & 973,445 & 3.2 & 3.2 & 1.0 & 1.9 & 5.5 & 19.4 \\
\bottomrule
\end{tabular}}
\begin{tablenotes}
\small \item Firm-year panel, 2017–2024. GHG Direct Intensity is P99-winsorised. Revenue in USD million.
\end{tablenotes}
\end{table}

Figure~\ref{fig:ghg_trends} plots mean GHG Direct Intensity over time, separately for treated and control groups under each treatment definition. The left panel shows the trajectory for firms in gas-exposed versus non-exposed countries; the right panel shows energy-intensive versus non-energy-intensive sectors. Both panels suggest a general declining trend in GHG intensity throughout the sample period, with a notable acceleration after 2022 for the energy-intensive sector group — which is the pattern our DiD specification is designed to identify.

\begin{figure}[H]
    \centering
    \includegraphics[width=\textwidth]{final_output/ghg_trends_plot.png}
    \caption{Mean GHG Direct Intensity over time, by treatment group. Left: gas-exposed vs.\ control countries. Right: energy-intensive (EII) vs.\ control sectors. Dotted vertical line marks the Russia–Ukraine shock (February 2022).}
    \label{fig:ghg_trends}
\end{figure}

\subsection{Treatment Assignment 1: Country-Level Russian Gas Exposure}

The first treatment dimension captures the degree to which a firm's host country was structurally dependent on Russian natural gas prior to the 2022 shock. We compute an annual country-level \textit{exposure ratio} defined as:
\[
  \text{Exposure}_{c,t} = \frac{\text{Russian gas imports}_{c,t}}{\text{Gross Inland Consumption}_{c,t}}
\]
using Eurostat datasets \texttt{nrg\_ti\_gas} (natural gas imports by partner country, in TJ gross calorific value) and \texttt{nrg\_bal\_c} (complete energy balances, Gross Inland Consumption in TJ). Both series are at annual frequency and cover EU-27 countries from 2015 to 2024.

To define the treatment group, we compute each country's average pre-shock exposure over the window 2017–2021 and classify countries in the top quartile ($\geq$ p75 = 14.1\%) as treated. This threshold yields eight treated countries: Hungary, Latvia, Slovakia, Italy, Czech Republic, Germany, Lithuania, and Bulgaria — precisely the group that features most prominently in policy discussions about Russian gas dependency. Figure~\ref{fig:exposure_bar} shows the full distribution of pre-shock average exposure across EU countries and the treatment threshold.

\begin{figure}[H]
    \centering
    \includegraphics[width=0.85\textwidth]{final_output/russian_exposure_barplot.png}
    \caption{Average Russian gas exposure by country, pre-shock window 2017–2021. Exposure is Russian gas imports as a share of gross inland energy consumption (Eurostat). Treated group (red): countries above the 75th percentile. Dashed lines mark the p75 and p50 thresholds.}
    \label{fig:exposure_bar}
\end{figure}

Figure~\ref{fig:exposure_trends} further illustrates the relevance of this classification by plotting the average gas exposure over time for treated and control countries. The divergence between the two groups is stark: treated countries had systematically higher exposure throughout the pre-shock period, and their exposure dropped sharply only after 2022 — precisely as Russian pipeline flows were curtailed.

\begin{figure}[H]
    \centering
    \includegraphics[width=0.75\textwidth]{final_output/exposure_trends_plot.png}
    \caption{Average Russian gas exposure over time for treated vs.\ control countries. Treated = top-p75 pre-shock exposure group (Hungary, Latvia, Slovakia, Italy, Czech Republic, Germany, Lithuania, Bulgaria). Dotted line: February 2022 shock.}
    \label{fig:exposure_trends}
\end{figure}

\subsection{Treatment Assignment 2: Energy-Intensive Sectors}

The second treatment dimension captures whether a firm belongs to a sector that is structurally energy-intensive. We classify sectors following the official Energy-Intensive Industries (EII) definition used by the OECD (based on energy input cost shares from the ICIO input-output tables), the European Commission (the EU EII Industrial Ecosystem), and the IEA (high-temperature process classification). Specifically, the following NACE Rev.2 two-digit codes are assigned $\text{EII}=1$:

\begin{center}
\small
\begin{tabular}{ll}
\toprule
NACE & Sector \\
\midrule
C16 & Manufacture of wood and wood products \\
C17 & Manufacture of paper and paper products \\
C19 & Manufacture of coke and refined petroleum products \\
C20 & Manufacture of chemicals and chemical products \\
C22 & Manufacture of rubber and plastic products \\
C23 & Manufacture of other non-metallic mineral products (cement, glass, ceramics) \\
C24 & Manufacture of basic metals \\
\bottomrule
\end{tabular}
\end{center}

The control group ($\text{EII}=0$) consists of structurally comparable manufacturing sectors that are not classified as energy-intensive: pharmaceuticals (C21), machinery and equipment (C28), motor vehicles (C29), fabricated metal products (C25), furniture (C31), and printing and media reproduction (C18). These sectors share key structural features with EII sectors — capital intensity, trade exposure, manufacturing scale — but have substantially lower energy input cost shares, making them suitable counterfactuals. In the final panel, 647,026 observations (66.4\%) belong to EII sectors and 327,299 (33.6\%) to the control group.

% =============================================================================
\section{Methodology}
\label{sec:method}
% =============================================================================

\subsection{Difference-in-Differences Framework}

We estimate the effect of the Russia–Ukraine energy shock on firm-level GHG Direct Intensity using a difference-in-differences (DiD) design. The identifying variation comes from comparing the change in GHG intensity for treated firms (those exposed to the shock through the two treatment dimensions defined above) relative to control firms, across the pre-shock (2017–2021) and post-shock (2022–2024) periods.

The baseline specification takes the form:
\begin{equation}
Y_{i,t} = \alpha_i + \lambda_t + \beta \cdot \text{Post}_t \times D_i + \gamma' X_{c,t} + \varepsilon_{i,t}
\label{eq:did}
\end{equation}
where $Y_{i,t}$ is GHG Direct Intensity for firm $i$ in year $t$; $\alpha_i$ are firm fixed effects, absorbing all time-invariant firm characteristics; $\lambda_t$ are year fixed effects, absorbing aggregate time trends; $D_i$ is one of the treatment indicators defined above; $\text{Post}_t = \mathbf{1}(t \geq 2022)$; and $X_{c,t}$ includes time-varying country-level controls — GDP per capita (in EUR) and the annual consumer price inflation rate. Standard errors are clustered at the country level to account for within-country correlation across firms and years.

We estimate four specifications. Models (1) through (3) each use a single binary treatment: country-level gas exposure (\texttt{gas\_exposed}), sector energy intensity (\texttt{energy\_intensive}), and the intersection of both (\texttt{both\_treated}), respectively. Model (4) — our most informative specification — includes all three treatment-post interactions simultaneously:
\begin{equation}
Y_{i,t} = \alpha_i + \lambda_t
  + \beta_1 \, \text{Post}_t \times \text{GasExp}_i
  + \beta_2 \, \text{Post}_t \times \text{EII}_i
  + \beta_3 \, \text{Post}_t \times \text{GasExp}_i \times \text{EII}_i
  + \gamma' X_{c,t} + \varepsilon_{i,t}
\label{eq:triple}
\end{equation}
This triple-interaction specification allows us to decompose the country-level and sector-level channels and estimate whether the two dimensions of treatment reinforce or partially offset each other.

\subsection{Identification: Exogeneity of the Shock}

A central advantage of this design is the exogeneity of the treatment timing. The Russia–Ukraine war and the associated energy disruptions were not driven by firm-level decisions: they represent a geopolitical shock that was unanticipated by firms in its timing and severity. This is in contrast to, for example, a carbon tax, which firms may anticipate and respond to in advance of implementation. Under this assumption, the 2022 shock constitutes a natural experiment, and the DiD estimator identifies the causal effect of exposure to the shock on GHG intensity.

\subsection{Parallel Trends Assumption}

The validity of the DiD estimate rests on the parallel trends assumption: in the absence of the treatment, treated and control firms would have followed similar trends in GHG intensity. We provide two pieces of evidence in support of this assumption. First, Figure~\ref{fig:ghg_trends} shows that treated and control firms exhibited broadly parallel declining trends in GHG intensity over the pre-shock period 2017–2021, with no systematic pre-shock divergence. Second, we estimate an event-study version of equation~(\ref{eq:did}), replacing $\text{Post}_t \times D_i$ with a full set of year-by-treatment interactions, normalised to zero in 2021 (the last pre-shock year). The coefficients on the pre-2022 interactions allow a direct test of pre-trend parallelism.

\begin{figure}[H]
    \centering
    \includegraphics[width=0.85\textwidth]{final_output/event_study_plot.png}
    \caption{Event study estimates: year-by-year coefficients for gas-exposed vs.\ control firms, relative to 2021. Pre-2022 coefficients test the parallel trends assumption. Shaded areas: 95\% confidence intervals. Firm and year FE; SE clustered by country.}
    \label{fig:event_study}
\end{figure}

The event study in Figure~\ref{fig:event_study} shows that pre-2022 coefficients are small and broadly centred around zero, consistent with parallel pre-trends. A post-2022 divergence becomes visible, though the pattern in the gas-exposure dimension is more modest than in the sector-level specification (see Figure~\ref{fig:ghg_trends}, right panel), which is consistent with the triple-interaction results discussed below.

% =============================================================================
\section{Results and Discussion}
\label{sec:results}
% =============================================================================

Table~\ref{tab:did} presents the four DiD estimates.

\begin{table}[H]
\centering
\small
\caption{Difference-in-Differences: Effect on GHG Direct Intensity}
\label{tab:did}
\begin{tabular}{lcccc}
\toprule
 & (1) & (2) & (3) & (4) \\
 & Gas Exposed & Energy Intensive & Both Treated & Triple Int. \\
\midrule
Post $\times$ Gas Exposed         & 5.206$^{***}$ &               &                & $-$2.602 \\
                                   & (1.100)        &               &                & (2.492) \\[2pt]
Post $\times$ Energy Intensive    &               & $-$38.468$^{***}$ &            & $-$43.477$^{***}$ \\
                                   &               & (2.210)           &            & (1.674) \\[2pt]
Post $\times$ Both Treated        &               &               & $-$14.582$^{***}$ & 9.846$^{***}$ \\
                                   &               &               & (5.601)           & (2.554) \\[2pt]
GDP per capita (EUR)               & 0.000$^{***}$ & 0.000$^{**}$  & 0.000          & 0.000$^{***}$ \\
                                   & (0.000)       & (0.000)       & (0.000)        & (0.000) \\[2pt]
Inflation (\%)                    & 0.219         & 0.395$^{**}$  & 0.802$^{***}$  & 0.239 \\
                                   & (0.205)       & (0.176)       & (0.300)        & (0.244) \\[2pt]
\midrule
Observations   & 973,445 & 973,445 & 973,445 & 973,445 \\
Within $R^2$   & 0.011   & 0.047   & 0.015   & 0.048 \\
Firm FE        & Yes     & Yes     & Yes     & Yes \\
Year FE        & Yes     & Yes     & Yes     & Yes \\
SE clustered at & Country & Country & Country & Country \\
\bottomrule
\end{tabular}
\begin{tablenotes}
\small
\item Notes: Dependent variable: GHG Direct Intensity (tCO$_2$e per \$1M revenue). Firm and year fixed effects included in all specifications. Standard errors in parentheses, clustered at the country level. $^{*}p<0.10$, $^{**}p<0.05$, $^{***}p<0.01$.
\end{tablenotes}
\end{table}

The results tell a clear and internally consistent story. Model (2) establishes the central finding: energy-intensive firms reduced their GHG Direct Intensity by approximately 38.5 tCO$_2$e per \$1M of revenue in the post-2022 period, relative to non-EII firms in the same countries and year ($p < 0.01$). This is a large and economically significant effect — roughly 32 percent of the sample mean intensity of 121 tCO$_2$e/\$M. The most natural interpretation is that the energy shock forced firms in high-energy-cost sectors to rethink their energy sourcing strategy: investing in energy efficiency, shifting to lower-carbon electricity contracts, or diversifying supply chains away from gas. Because natural gas is itself a GHG-emitting fuel, this diversification toward electricity (especially from renewable or nuclear sources) mechanically reduces Scope 1 emissions per unit of output.

Model (4) — the triple-interaction specification — allows us to decompose the mechanism further. The coefficient on Post $\times$ EII ($\hat\beta_2 = -43.5$, $p<0.01$) remains large and highly significant. The coefficient on Post $\times$ Both Treated ($\hat\beta_3 = +9.8$, $p<0.01$) is positive and significant, indicating that the energy-transition effect is \textit{attenuated} for energy-intensive firms located in countries that were more dependent on Russian gas. The net effect for a firm that is both in a gas-dependent country and in an EII sector is therefore $-43.5 + 9.8 = -33.7$ tCO$_2$e/\$M — substantial, but meaningfully smaller than the broader EII-sector average.

The attenuation effect is theoretically compelling. Countries with high pre-shock Russian gas exposure faced a more constrained energy transition path: the alternatives to pipeline gas — LNG imports, interconnectors, domestic renewable expansion — are subject to infrastructure bottlenecks and long lead times. The contrast between France, which could accelerate its nuclear fleet utilisation, and Italy, which at various points considered reopening mothballed coal capacity, illustrates the mechanism. Firms headquartered in more constrained energy systems had fewer viable low-carbon alternatives in the short run, limiting the GHG reductions they could achieve through energy diversification. This finding suggests that firm-level decarbonisation and country-level energy policy are deeply complementary: even when firms face strong incentives to reduce emissions, their ability to act is constrained by the energy infrastructure available to them.

Model (1) — the pure country-exposure specification — shows a small but significantly positive interaction coefficient ($\hat\beta_1 = +5.2$, $p<0.01$). Taken in isolation, this seems to suggest that firms in gas-exposed countries increased their GHG intensity. However, Model (4) reveals that this is largely driven by the composition of sectors across countries: once we control for the EII/non-EII split, the standalone gas-exposure effect becomes statistically indistinguishable from zero ($-2.6$, $p=0.30$). The apparent positive association in Model (1) reflects the fact that gas-dependent countries host proportionally more EII firms, and those firms' transition costs dominate the country-level average.

% =============================================================================
\section{Limitations and Further Research}
\label{sec:limits}
% =============================================================================

Several limitations of the current analysis are worth acknowledging. First, and most importantly, the country assignment is based on each firm's registered headquarters location. For MNEs, this may poorly capture the true geography of production and energy consumption. A firm registered in Germany but operating primarily in Poland faces the energy mix of Poland, not Germany — and Germany's high gas exposure score would be attributed to a firm that may not bear that exposure operationally. Correcting for this would require facility-level data on production locations, which is not available in TruCost.

Second, the analysis cannot distinguish between genuine decarbonisation and carbon leakage. If firms in gas-exposed countries reduced their measured Scope 1 emissions by outsourcing the most emission-intensive production steps to lower-cost jurisdictions outside Europe, our estimate would overstate the true environmental impact of the energy shock. A Scope 3 analysis or trade-adjusted emission accounting would be needed to address this.

Third, the parallel trends assumption, while supported by the event study, cannot be formally proven. The pre-shock period 2017–2021 contains other shocks (notably COVID-19 in 2020) that may have differentially affected treated and control firms. Although we control for GDP per capita and inflation, the COVID-19 recession and recovery pattern could introduce confounding variation that is not fully absorbed by the year fixed effects.

A natural extension of this work is to examine the medium-term trajectory. The current panel extends to 2024, capturing only two full years of the post-shock period. It is plausible that the transition away from Russian gas is a multi-year process, with the full decarbonisation effect materialising over a longer horizon as firms complete capacity investments in renewables and energy efficiency. A repeated analysis with data from 2025 and beyond would allow for a richer dynamic specification.

% =============================================================================
\section{Conclusion}
\label{sec:conclusion}
% =============================================================================

This paper provides firm-level evidence that the Russia–Ukraine energy shock triggered a significant reduction in GHG intensity among European energy-intensive firms. Using a difference-in-differences design on a panel of approximately 973,000 firm-year observations, we find that energy-intensive firms reduced their direct GHG emissions per unit of revenue by approximately 38--43 tCO$_2$e per \$1M in the post-2022 period, relative to comparable non-EII firms. The most plausible mechanism is that the sharp increase in energy prices — and the strategic imperative to diversify away from Russian gas — pushed firms to invest in energy efficiency and lower-carbon energy sources, reducing Scope 1 intensity as a by-product.

Critically, this effect is not uniform: firms headquartered in countries with higher pre-shock Russian gas dependency experienced an attenuation of roughly 10 tCO$_2$e/\$M, leaving them with a smaller but still substantial net reduction. This finding highlights the strategic interdependence between firm-level environmental performance and country-level energy infrastructure. Countries with more diversified, flexible energy systems — nuclear, renewables, well-connected grids — were better positioned to channel the energy price shock into genuine decarbonisation gains. Countries locked into pipeline gas dependency faced a slower transition, constraining the environmental improvements their industrial firms could achieve even when strong price incentives were present.

The policy implication is dual. At the country level, investment in energy diversification — transmission interconnectors, LNG import capacity, domestic renewable deployment — is not only an energy security imperative but also an enabler of corporate decarbonisation. At the firm level, companies constrained by their location should place strategic emphasis on energy-mix flexibility and long-term power purchase agreements as tools to decouple their GHG performance from the energy endowment of their host country.

% =============================================================================
\section*{References}
% =============================================================================

\setlength{\parindent}{0pt}

Abdullah, M., Tiwari, A.K., Hossain, M.R. and Abakah, E.J.A. (2024) `Geopolitical risk and firm-level environmental, social and governance (ESG) performance', \textit{Journal of Environmental Management}, 363, 121245.

Barkemeyer, R., Revelli, C. and Douaud, A. (2023) `Selection bias in ESG controversies as a risk for sustainable investors', \textit{Journal of Cleaner Production}, 405, 137035.

Bennett, N. and Lemoine, G.J. (2014) `What a difference a word makes: Understanding threats to performance in a VUCA world', \textit{Business Horizons}, 57(3), pp.\ 311--317.

Bolton, P. and Kacperczyk, M. (2023) `Global Pricing of Carbon-Transition Risk', \textit{The Journal of Finance}, 78(6), pp.\ 3677--3754.

Card, D. and Krueger, A.B. (1994) `Minimum Wages and Employment: A Case Study of the Fast-Food Industry in New Jersey and Pennsylvania', \textit{American Economic Review}, 84(4), pp.\ 772--793.

Cavusgil, S.T. et al.\ (2021) `International business in an accelerated VUCA world: Trends, disruptions, and coping strategies', \textit{Rutgers Business Review}, 6(3), pp.\ 219--243.

Collste, D. et al.\ (2025) `Polycrisis patterns: applying system archetypes to crisis interactions', \textit{Global Sustainability}, 8, e17, pp.\ 1--11.

Erzurumlu, Y.O. et al.\ (2025) `The effects of geopolitical and political risks on corporate ESG practices', \textit{Journal of Environmental Management}, 384, 125747.

Fan, D., Zhou, Y., Yeung, A.C.L., Lo, C.K.Y. and Tang, C.S. (2022) ‘Impact of the U.S.–China trade war on the operating performance of U.S. firms: The role of outsourcing and supply base complexity’, Journal of Operations Management, 68(8)

Feng, X. et al.\ (2025) `Russia--Ukraine conflict and local government debt risk in China', Working paper.

Future Earth, Sustainability in the Digital Age, and International Science Council (2021) \textit{Our Future on Earth}.

Gereffi, G. (2014) `Global value chains in a post-Washington Consensus world', \textit{Review of International Political Economy}, 21(1).

Gereffi, G., Humphrey, J. and Sturgeon, T.J. (2005) `The governance of global value chains', \textit{Review of International Political Economy}, 12.

Henderson, J. et al.\ (2002) `Global production networks and the analysis of economic development', \textit{Review of International Political Economy}, 9.

Homer-Dixon, T. et al.\ (2015) `Synchronous failure: The emerging causal architecture of global crisis', \textit{Ecology and Society}, 20(3), art6.

Hover, R., Buchenrieder, G. and Sokolsky, J. (2025) `Economic interdependence and the likelihood of war -- a systematic literature review', \textit{Cogent Social Sciences}, 11(1).

Humphrey, J. and Schmitz, H. (2002) `How does insertion in global value chains affect upgrading in industrial clusters?', \textit{Regional Studies}, 36(9).

Kano, L. (2018) `Global value chain governance: A relational perspective', \textit{Journal of International Business Studies}.

Lawrence, M., Janzwood, S. and Homer-Dixon, T. (2022) `What Is a Global Polycrisis?', Version 2.0, Discussion Paper 2022-4. Cascade Institute.

Lessard, D., Teece, D.J. and Leih, S. (2016) `The dynamic capabilities of meta-multinationals', \textit{Global Strategy Journal}, 6(3), pp.\ 211--224.

Mansfield, E. and Pollins, B. (2003) `Interdependence and conflict: An Introduction', in Mansfield, E. and Pollins, B. (eds.) \textit{Economic Interdependence and International Conflict}. University of Michigan Press.

McWilliam, S.E. et al.\ (2020) `Global value chain governance: Intersections with international business', \textit{Journal of World Business}, 55.

Network for Greening the Financial System (NGFS) (2024) \textit{Improving Greenhouse Gas Emissions Data}.

Olden, A. and Moen, J. (2022) `The Triple Difference Estimator', \textit{Econometrics Journal}, 25(3), pp.\ 531--553.

Perdichizzi, S. et al.\ (2024) `Carbon emission and firms' value: Evidence from Europe', \textit{Energy Economics}, 131, 107324.

Petricevic, O. and Teece, D.J. (2019) `The structural reshaping of globalization: Implications for strategic sectors, profiting from innovation, and the multinational enterprise', \textit{Journal of International Business Studies}, 50, pp.\ 1487--1512.

Ponte, S. (2020) `The hidden costs of environmental upgrading in global value chains', \textit{Review of International Political Economy}.

Rapoza, K. (2019) `The unintended consequences of the China trade war', \textit{Forbes}, June 26.

Renn, O. et al.\ (2020) `Systemic Risks from Different Perspectives', \textit{Risk Analysis}.

Renn, O. et al.\ (2019) `Things Are Different Today: The Challenge of Global Systemic Risks', \textit{Journal of Risk Research}, 22(4), pp.\ 401--415.

Robertson, R. (1983) `Interpreting Globality', in \textit{World Realities and International Studies Today}. Glenside, PA: Pennsylvania Council on International Education.

Rugman, A. and Oh, C.H. (2008) `Friedman's follies: Insights on the globalization/regionalization debate', \textit{Business and Politics}, 10(2), pp.\ 1--14.

Saharti, M. et al.\ (2024) `Environmental, social and governance (ESG) performance of firms in the era of geopolitical conflicts', \textit{Journal of Environmental Management}, 351, 119744.

Sautner, Z. et al.\ (2023) `Firm-Level Climate Change Exposure', \textit{The Journal of Finance}, 78(3), pp.\ 1449--1498.

Schmidt-Felzmann, A. (2016) `It's never ``just business''', \textit{Global Affairs}, 2(2), pp.\ 115--118.

Schoemaker, P.J.H., Heaton, S. and Teece, D.J. (2018) `Innovation, dynamic capabilities, and leadership', \textit{California Management Review}, 61(1), pp.\ 15--42.

S\&P Global Sustainable1 (2025) \textit{Trucost Environmental Data: Methodology}.

Schweizer, P.-J. and Renn, O. (2019) `Governance of Systemic Risks for Disaster Prevention and Mitigation', \textit{Disaster Prevention and Management}, 28(6), pp.\ 862--874.

U.S. Environmental Protection Agency (EPA) (2025) \textit{Scope 1 and Scope 2 Inventory Guidance}.

van Tulder, R., Verbeke, A. and Jankowska, B. (eds.) (2020) \textit{International Business in a VUCA World}, PIBR Vol.\ 14. Bingley: Emerald Publishing.

Yeung, H.W.-c. and Coe, N.M. (2015) `Toward a Dynamic Theory of Global Production Networks', \textit{Economic Geography}, 91(1), pp.\ 29--58.

\end{document}
